\documentclass[output=paper]{langscibook}
\ChapterDOI{10.5281/zenodo.18845920}
\author{Peredur Webb-Davies\orcid{}\affiliation{Bangor University} and María del Carmen Parafita Couto\orcid{}\affiliation{University of Santiago de Compostela} and Draško Kašćelan\orcid{}\affiliation{University of Essex}}
\title{Code-switching in context: Variation, processing, and development}
\abstract{\noabstract}
\IfFileExists{../localcommands.tex}{
  \addbibresource{../localbibliography.bib}
  \usepackage{tabularx,multicol}
\usepackage{url}
\urlstyle{same}

 
\usepackage{langsci-optional}
\usepackage{langsci-lgr}
\usepackage{langsci-gb4e}

\usepackage{multirow}
\usepackage{cleveref}
\usepackage{subfigure}
\usepackage{langsci-branding}

  \input{../localcommands} 
  \input{../localhyphenation} 
  \togglepaper[1]%%chapternumber
}{}

\begin{document}

\maketitle 
%\shorttitlerunninghead{}%%use this for an abridged title in the page headers

\section{Volume introduction}

This volume provides an exploration of code-switching in bilinguals{\slash}multilin\-guals, examining the phenomenon in linguistic development, social interaction, and processing. The volume focuses on grammatical variation in speech language patterns, as well as the development of code-switching, including its occurrence in neurodevelopmental conditions. By integrating research across diverse bilingual contexts, this collection offers a comprehensive understanding of how code-switching functions within different linguistic communities.

Code-switching (or code-mixing) is a well-documented feature of the speech of multilinguals whereby words or grammar from multiple languages are combined in a single discourse, conversation or utterance. Use of code-switching is shaped by a confluence of factors, including cognitive development, linguistic structures and, crucially, the socio-cultural contexts in which bilinguals live and interact. The chapters in this volume investigate code-switching from different perspectives and explore other governing factors at play.

In grammatical terms, code-switching can be classified in various ways, perhaps the most salient being the distinction between inter-clausal switching (where the switch occurs at clause boundaries) and intra-clausal switching (where the switch happens within a clause).\footnote{Note that across literature, there is a variation in whether a clause or a sentence is used as a unit of analysis. In case a sentence is used as a unit of analysis, intra-sentential switch would refer to a switch within a sentence (which could occur between clauses in case of complex or compound sentences), while inter-sentential switch would refer to a switch between sentences.} Research on code-switching has been prevalent in the literature for many decades, and since the 1980s a particular focus has been given to grammatical constraints (if any) on code-switching. A brief reference to the journey of this debate is given by Valdés Kroff et al. in chapter 5 of this volume. It will suffice to note at present that studies in this field have investigated the patterns of code-switching found in many diverse speech populations, with a particular focus since the work of \citet{Poplack1980} on Spanish-English, but also a significant contribution to the study of the speech of Welsh-English bilinguals by Deuchar and colleagues (e.g., \citealt{DeucharEtAl2014, DeucharEtAl2018}). Generally, studies show that code-switching is not the same for all bilinguals (cf. the in-depth meta-review and typology provided by \citealt{Muysken2001}), though it is by no means random and is in fact governed by factors such as community norms and the mapping of the grammars of the languages involved.

Chapters 2--8 in this volume comprise a range of studies examining code-switching from multiple perspectives. Key themes explored across the presented studies include:

\begin{itemize}
\item \emph{Community norms and language practices:} Code-switching reflects the linguistic norms of a given community. When bilingualism is widely accepted and practiced, children are more likely to integrate code-switching into their speech. Conversely, in communities where strict language separation is emphasized, code-switching may be less frequent or carry a different social meaning (chapter 6). Code-switching for heritage communities (such as the Piedmontese-Spanish speaking communities in Argentina) can be limited to certain kinds of activities where the heritage language is prioritised, though switching still occurs (chapter 2). From a methodological standpoint, tasks that are minimally metalinguistic -- such as a participants indicating they comprehend a stimulus -- appear to be more in-line with the community speech norms of bilinguals, whereas more metalinguistic tasks provide different results (chapter 5).
\item \emph{Caregiver influence on bilingual development:} Parents and caregivers play a pivotal role in shaping children's code-switching behaviors. Their choices in providing bilingual input, whether through inter-sentential or intra-sentential switching, impact children's language development. Studies suggest that children tend to accommodate their caregivers in speech practices, including the use of code-switching (chapter 7).
\item \emph{Attitudes and ideologies surrounding bilingualism:} Perceptions of code-switching vary widely. In some contexts, bilingualism is celebrated, and code-switching is viewed as a resourceful linguistic tool. In others, it is stigmatized as a sign of incomplete language acquisition. These attitudes shape language policies at family, school, and community levels, affecting bilingual children's linguistic experiences (chapter 8). Bilinguals recognize that code-switching within phrases (e.g., determiner phrases) occur in their community’s speech, but they are found in a study in this volume to judge them less grammatically acceptable than unilingual utterances from one of their languages (chapter 4).
\item \emph{Code-switching as a communicative resource:} Rather than being a linguistic deficit, code-switching serves as a strategic tool for meaning-making and social identity negotiation. Bilingual children use code-switching to enhance their expressive capabilities and engage more effectively in communication (chapter 7). Young bilinguals who use instant messaging deploy code-switching creatively in terms of both word-selection, orthography and degree of morphological integration of switched items (chapter 3).
\item \emph{Neurodiversity and bilingualism:} By reviewing research on bilinguals with neurodevelopmental conditions, this volume broadens the understanding of code-switching. It challenges deficit-based perspectives that pathologize language mixing and highlights the importance of evidence-based guidance for multilingual families with neurodivergent family members (chapter 8).
\end{itemize}

The remainder of this chapter summarizes and contextualizes the studies presented in this book.

% \section{\sffamily}
% \subsection{\sectref{sec:key:1}: Grammatical variation and code-switching}
\subsection{Grammatical variation and code-switching}

The range of approaches to researching code-switching variation in the field is well-represented in chapters 2--5, with brief commentaries on those chapters provided here.

The study by Goria and Bellamy (chapter 2) investigates code-mixing in the speech of Heritage Piedmontese-Argentinian Spanish bilinguals living in Argentina, specifically focusing on Spanish noun insertions into Piedmontese matrices. Note that the authors of the study distinguish code-mixing from code-switching on the basis of pragmatic function. This chapter is noteworthy as the first corpus-based study into code-mixing in Heritage Piedmontese.

The authors draw their data (a sample of 22 speakers) from a larger corpus of spontaneous speech collected as part of an ongoing project focusing on the linguistic variety Heritage Piedmontese (HP). The corpus data includes semi-structured interviews between a native HP researcher and one or more HP speakers, with the group conversations involving more spontaneous participant-to-participant interaction. The authors highlight that the active use of HP among the community was primarily found during specifically Piedmontese-related cultural activities, and the recorded conversations took place in an explicit monolingual mode; that is, participants were requested to use HP, though this did not preclude code-mixing of the type which is described by the authors.

They identify several distinct types of switching of Spanish noun material in their data, with insertion being much more common than alternation, although the authors note that this preference is not found to such an extent in the other data they collected (i.e., when not in an interview setting). The most frequently-observed insertional type in their data is what they call “minimal insertions”, i.e., the insertion of Spanish nouns either as bare forms (which comprises the majority of the tokens in their analysis, 64\%) or with morphological Piedmontese morphology which serves to integrate switched words with the Piedmontese grammar used in the rest of the sentence. A significant difference was found for grammatical number, where plural nouns were most likely be inserted with Piedmontese morphological marking than singular nouns, and a further significant difference was observed that contrasts inserted feminine Spanish nouns, which are more likely to be integrated to HP morphology, than masculine nouns, which are more likely to be bare forms.

The authors explain the patterns found as being emblematic of a preference by these participants to “display their knowledge” of Piedmontese and thereby avoid code-mixing where possible, a practice observed in other studies of similar social contexts (e.g., \citet{CarterEtAl2011} for Welsh-Spanish bilinguals in Argentina). Regarding the patterns shown with noun number and gender, the authors interpret this to relate to the grammatical congruence between the two languages (both Romance), i.e., at points of congruence between HP and Spanish, the same grammatical structure is maintained. A small number of instances of insertions are found which were not predicted by the authors, and which they suggest may indicate innovative forms.

The interaction between languages explored in this study can be classified as between a nationally prestigious language (Argentinian Spanish, the primary language of that country) and a heritage language (Heritage Piedmontese, originating in Piedmont, Italy). Piedmontese is endangered in Piedmont and has restricted usage in Argentina as the language of heritage communities, though it has experienced a degree of language revival. Literature on situations involving speakers who are bilingual in a majority and a minority language has demonstrated the common influence of the former on the latter, such as in the form of language shift – characterized as the diachronic preference of speakers in a speech community for one language over another less prestigious one (e.g., \citealt{AppelMuysken2006}) – or linguistic interference (e.g., \citealt{ThomasonKaufman1988}). Case studies of the influence of code-switching in particular on grammatical changes within the minority language in such cases include the work of Deuchar and colleagues on Welsh-English bilinguals (e.g., \citealt{DeucharDavies2009}). The study presented in this chapter provides a valuable addition to this field.

Additionally, the use of recorded spontaneous speech to analyse language use patterns is, of course, well-established in sociolinguistics. Numerous past projects have collated corpora of bilinguals’ speech in a similar fashion, including, notably, the corpora of Welsh-English, Spanish-English and Welsh-Spanish corpora (the latter likewise from Argentina) by researchers based in Bangor University, Wales, under the auspices of Margaret Deuchar (q.v. \citealt{DeucharEtAl2018}). That data is publicly accessible (\citealt{Bangortalknodate}), and indeed we view Margaret Deuchar as a pioneer in making bilingual linguistic data widely available for use. It is hoped that much more bilingual corpora are made openly accessible for researchers to verify and explore, especially taking into account considerable hard work involved in collecting, coding and transcribing corpus data in the first place (see, e.g., the Talkbank repository of corpus data under the stewardship of Brian \citealt{MacWhinney2000}); indeed, it is commendable that Goria and Bellamy provided a link to their data via the Open Science Framework website.

Another corpus study, by Vihman and colleagues, is presented in chapter 3, with related data being made available on the Open Science Framework. The authors investigate code-switching in instant messaging (IM) by Estonian-English bilinguals aged under 18, focusing on the extent and nature of the morphological and orthographic integration of English elements into an Estonian frame. The teenagers under study live offline in a relatively monolingual Estonian sociolinguistic context, but engage with online chat in a much more multilingual manner \citep{VihmanEtAl2022}, using English particularly for stylistic or pragmatic function within their linguistic repertoire. The authors look at the test case of the language used in instant messaging sourced from a corpus that consists of IM from over 100 Estonian teenage participants; the study presented here includes data from 99 participants. A clear contribution of this chapter is therefore its exploration of code-switching in computer-mediated communication, a mode that is text-based, and therefore governed to some extent by the orthographic conventions of the language in question, but is also frequently influenced by features of the spoken informal language, including not only phonological features but also code-switching, as the authors point out.

In general, Vihman et al. find a low proportion of English words in the dataset (7.5\% of tokens; in contrast, speech data analyzed by \citet{VihmanEtAl2022} had 3\% English tokens), and it is observed that these participants generally use an Estonian morphosyntactic frame for their IMs. The authors present a detailed quantitative analysis of orthographic and morphological treatment of English insertions, including examining participant variables and different parts of speech. Orthographic integration of English insertions into the Estonian frame was low: only 17\% of tokens were either orthographically integrated or had a mixture of elements of English and Estonian orthography; the remainder were coded as English switches which retained English spelling. The results show some variation between participants based on factors including age, with younger participants integrating less than older ones. With regard to morphological integration, the authors note that most switched material in their data is not able to be integrated anyway, such as interjections and adverbs, but of the switches where morphological integration was available, the majority (79\%) are indeed integrated, in contrast to the orthographic data, with less inter-participant variation. There are very few tokens in the dataset which exhibit both orthographic and morphological integration. The authors provide a detailed discussion of the interplay between orthographic and morphological treatment of English switching in the IM language of these teens, making a persuasive example of code-switching used for stylistic as well as communicative purposes.

The extent to which code-switched material from a language is integrated into the grammar of a recipient language has long been an interest of the field. Although some theoretical approaches argue that, e.g., the affixation of bound morphemes from one language onto a stem from another violates some constraint (cf. \citealt{Poplack1980}), and indeed some studies find bare forms (e.g., \citealt{Schmitt2000}); it has been demonstrated elsewhere that this occurs across different languages (e.g., \citealt{DeucharEtAl2018, Myers-Scotton2002}); see also Goria and Bellamy in chapter 2 of this volume. Indeed, the Matrix Language Frame model (\citealt{Myers-Scotton2002, Myers-ScottonJake2009}) posits that it is to be expected that an inserted lexical item is morphologically integrated into the morphosyntactic frame into which it is embedded, as found in, e.g., \citet{Deuchar2020} or \citet{DeucharEtAl2018} – assuming, of course, that the languages in question have morphology that enables this to be produced or observed. In the case of Estonian and English, Estonian is far more morphologically rich than English, having, for example, a complex system of inflectional case marking, and as such this provides an ideal context for identifying the extent of the integration of code-switched material. 

The findings in Vihman et al.’s study here show that, while IM language is broadly comparable to speech norms (at least in their case), there is added interest in examining it as it presents other factors, such as spelling, to consider, but also because in text discourse, as the authors demonstrate here, it is possible to evaluate more closely the apparent intent of users in how they are deploying code-switched material, e.g., via the extent to which they integrate the orthography and what this implies. We welcome the expansion of the study of code-switching into non-spoken contexts and look forward to more contributions like this.

In chapter 4, Ezeizabarrena, Munarriz-Ibarrola and deCastro-Arrazola describe their study of Basque-Spanish code-switching, focusing on the interaction of Basque and Spanish within determiner phrases (DP). So-called mixed DPs have been, as the authors note, the focus of several previous studies, including such language pairs as English-Spanish (\citealt{Dussias2003, ValdésKroff2016}) and Welsh-English (\citealt{ParafitaCoutoEtAl2013, ParafitaCoutoGullberg2019}); see \citealt{BellamyParafitaCouto2022} for an overview. In the case of languages with different gender rules – and, e.g., different ways of marking gender within DPs – this is a particularly interesting focus for analysing code-switching, since speakers producing such phrases must select some combination of morphemes from the two interacting languages, which may or may not adhere to the grammar of one or both those languages. In the case of Basque and Spanish, Spanish nouns have gender morphology whereas Basque does not; furthermore, there are other differences with the behaviour of elements such as determiners between the two languages. Previous studies have identified varying patterns of preference by judges on how Basque nouns interact with Spanish determiners (\citealt{BadiolaSande2018, IriondoEtxeberria2017, ParafitaCoutoEtAl2015}), which \citet{Munarriz-IbarrolaEtAl2021} demonstrate is linked to a bilingual’s first language. Ezeizabarrena and colleagues elicit grammaticality judgments from two groups of Basque-Spanish bilinguals, L1 Spanish (also comprising simultaneous bilinguals and early sequential bilinguals) and L1 Basque (all early sequential bilinguals). Their general aim was to identify which patterns of integration of Basque nouns in mixed DPs are preferred by the participants.

The method implemented in this study required participants to listen to stimuli sentences containing mixed Basque-Spanish DPs comprising a Spanish determiner preceding a Basque noun, and to rate their acceptability on a scale from 1 (bad) to 5 (good). The results show a mean acceptability of 3.54 for the critical sentences, which the authors interpret as general acceptance or familiarity with mixed DPs by these speakers, although participants did not rate such sentences as highly as unilingual Spanish or Basque sentences. The judgment patterns of both participant groups were broadly similar, although the L1 Spanish bilinguals’ responses were qualitatively more polarised than those of the L1 Basque bilinguals. The authors conduct statistical tests that identify some patterns that contrast the two groups’ responses, though no clear overall preference was shown to grammatical gender in mixed DPs, in contrast to previous studies. The authors conclude that their findings are indicative of speakers using word-specific strategies when it comes to gender marking in mixed DPs, or that certain kinds of bilingual (e.g., relating to language dominance) may prefer certain strategies.

Valdés Kroff et al.’s chapter 5 in this volume, following up on one of their  previous studies (\citealt{ValdésKroffEtAl2017}), report on their investigation into an established finding in the literature on Spanish-English code-switching, whereby switch points that occur before an auxiliary such as \textit{haber} ‘have’ are found much more commonly than switch points after such an auxiliary. Recruiting L1 English-L2 Spanish participants living in Florida, the authors aimed to test participants’ reactions to code-switched sentences of varying degrees of grammaticality (i.e., including both sentence types that are typically found in the speech of such bilinguals, as well as those that are seldom or never found). Spanish-English bilinguals in Florida have been the focus of numerous studies (e.g., \citealt{ParafitaCoutoEtAl2011,ParafitaCoutoEtAl2013}); speakers in this population have been argued to index code-switching with identity and to include both inter-sentential and intra-sentential switching in their repertoires (e.g., \citealt{CarterEtAl2011}). As a population that, varying by city or region, is intensely exposed to both Spanish and English – and Spanish-English code-switching – in their daily lives, such bilinguals are well-suited candidates for investigating the practice of and responses to code-switching from a psycholinguistic perspective.

In the present work by Valdés Kroff et al., participants undertook an eye-tracking experiment requiring them to read test sentences which included code-switching between Spanish and English, but which varied in terms of the switch points within the sentence: \textit{haber} ‘have’ + verb participle constructions with a switch either before or after \textit{haber}, and \textit{estar} ‘be’ + verb participle constructions with a switch either before or after \textit{estar}; sentences were also divided between perfect and present tense constructions. The participants were further contrasted based on what task they were asked to complete after reading each sentence: either to judge whether the sentence was grammatical or not, or to indicate whether or not they comprehended the sentence.

A key finding of this study is a difference between the group who did the comprehension task and those who did the grammaticality judgment task. Based on the early reading measure of gaze duration, the group participating in the comprehension task was found to read switches at the participle more slowly than at the auxiliary (with both auxiliary types -- \textit{haber} and \textit{estar}). However, based on the total duration, this group showed slower reading times at the participle than at the auxiliary when \textit{haber} was used, while there were no differences in relation to auxiliary \textit{estar}. On the other hand, the participant group who did the grammaticality judgment task did not demonstrate the same extent of variability in reading speed: for this group, overall trials on which switches occurred at the participle were read more slowly than at the auxiliary. This distinction between the two groups may be interpreted as evidence that being asked whether a sentence is grammatical is a metalinguistic task which represents natural language processing less well than when asked whether a sentence is meaningful, based on the evidence from the comprehension task group, whose responses are similar to the patterns already described in the literature.

The authors argue that the findings of their study point to the online processing of code-switching matching speech norms of the community; e.g., since a switching distinction between \textit{haber} and \textit{estar} is typical in speech among Spanish-English bilinguals in Florida (for example), that distinction is also reflected in the psycholinguistic processing by such bilinguals. The authors rightly point out that code-switching norms vary from community to community (q.v. \citealt{ParafitaCouto2017, CarterEtAl2011, Deuchar2020}), which is indeed motivation for further replication of their method to other contrastive speech communities, in order to see whether the findings would be repeated in populations where, for example, the languages in question or the code-switching habits of speakers are different.

This is a valuable study that further helps to bridge research in language from a social perspective and language from a psycholinguistic perspective: cf. earlier forays into this by \citet{Vaughan-EvansEtAl2020}, the authors of which were a mix of psychologists and linguists; that study also found, in broad terms, a parallel between electrophysiological evidence and the code-switching patterns found in the speech of Welsh-English bilinguals (e.g., \citealt{DeucharEtAl2018}).

% \subsection{\sectref{sec:key:2}: Development and neurodevelopmental conditions}
\subsection{Development and neurodevelopmental conditions}

Several chapters in this volume explore code-switching (predominantly) in bilingual children to examine the critical role of the surrounding linguistic environment. These chapters illuminate the ways in which community language, caregiver practices, and prevailing attitudes toward bilingualism converge to shape the code-switching patterns of young children, both neurotypical and neurodivergent.

Earlier approaches to code-switching have often focused on linguistic constraints, seeking to identify the grammatical principles that govern language mixing. While such analyses are valuable, they often overlook the dynamic interplay between the individual and their environment (for an overview see \citealt{Deuchar2020, ParafitaCoutoetal2023, PhillipsDeuchar2021}). The studies in this volume place code-switching within a broader ecological framework, recognizing that language is not simply a cognitive tool but a social practice embedded in specific communities and relationships. The ecological perspective adopted here highlights the interplay between individual linguistic development and broader socio-cultural contexts. 

In chapter 6, Vihman and Vihman review and analyse 47 studies on the spontaneous language use by bilingual children between the ages of 2--4 years. By looking at emergent grammars, they found that almost every aspect of early morphosyntax (apart from root infinitives) is open to the influence of bilingual interaction. In terms of type of interaction, they find evidence of acceleration, delay, and transfer between structures. The review remains inconclusive about the effect that dominance plays in interaction between structures, likely due to various ways in which dominance has been operationalised across studies. Vihman and Vihman point to the importance of individual differences in dynamic experiences of bilinguals, which might affect whether dominance comes up as an important factor. This review adds to the long standing discussions in the field of bilingualism about the lack of consensus on how specific phenomena, such as dominance, are defined (e.g., \citealt{Treffers-Daller2019}), which ultimately has an impact on comparisons between existing research and generalisations that can be drawn about bilingual development. Vihman and Vihman conclude that emerging bilingual grammars are moving targets -- consequently, it might be hard to estimate the required degree of overlap between languages to trigger interaction. As the previous overviews of bilingual grammars focused on experimental studies and included older children (e.g.,  \citealt{VanDijkEtAl2022}), Vihman and Vihman’s overview supplements their work by looking at spontaneous speech of emerging bilinguals. As they point out, small samples across identified 47 studies illustrate a need for large scale research within this age group. In light of that, their review is a vital state-of-the-art summary which maps out areas that require further attention, such as looking at patterns of use in the wider community and their effect on language use in young bilinguals.

Work by Quay (chapter 7) illustrates the value of an in-depth case study approach in understanding the dynamic nature of bilingual development in children, as well as the impact of caregivers, such as grandparents. The case study follows the linguistic development of Ryan, a Cantonese-English bilingual, between the ages of 2;9 and 4;11 in his interactions with his grandmother. Sixteen transcripts (totalling 460 minutes of video recordings) in which Ryan and his grandmother were alone were analysed to explore changes in his language use patterns. The study aimed to explore whether introducing more English (the societal language) through code-switching practices would facilitate Ryan’s frequency of use and skill development in English (his weaker language), and his overall development as a Cantonese-English bilingual in the Canadian context. Over the course of 2 years, data indicated that Ryan is generally accommodating his grandmother’s input in terms of the frequency patterns of using English-only, Cantonese-only, or code-switched utterances. In terms of intra-sentential code-switching (i.e., switching within a sentence), Ryan and his grandmother had a similar proportion of one-word insertions (62\% and 69\% respectively), two-word mixed utterances (8\% each), and multi-word mixed utterances (30\% and 23\% respectively). Further qualitative analyses identified that Ryan’s English utterances were morphosyntactically less complex than his Cantonese utterances, suggesting that despite accommodating his grandmother’s language choice, he remained a Cantonese-dominant bilingual. Unlike his grandmother’s input, Ryan’s speech showed evidence of congruent lexicalisation and convergence, which he used in more complex and in compound sentences, as well as to form complete utterances in his ‘weaker’ language. Quay argues that this data stresses the importance of analysing single-language utterances by emerging bilinguals, especially in their ‘weaker’ language, as they can offer insight into language mixing phenomena through convergence. This longitudinal observation shows how grandmother’s code-switching practices as well as the child’s pragmatic sensitivity to language choice facilitated the development of initially ‘weaker’ language (i.e., the societal language) of an emerging bilingual. Importantly, this data highlights that unrestricted use of both languages facilitates linguistic development of bilinguals as it allows drawing upon their complete linguistic repertoire. 

The scoping review by Kašćelan and Parafita Couto (chapter 8) outlines existing findings on code-switching by individuals with neurodevelopmental conditions. By looking at this under-explored area, they identify lack of representation in various aspects. Specifically, in their analysis of 31 manuscripts, the focus was primarily on code-switching in autism and language disorder, while code-switching in other neurodevelopmental conditions (e.g., ADHD, stuttering, etc.) has been rarely investigated. In terms of demographic-related data, most studies focused on children of age 8 or under, and with participants predominantly from North America and Europe. Code-switching in geographic locations where multilingualism and code-switching tend to be mainstream norms of communication are rarely explored. The most striking finding of this review points to the approaches in which code-switching has been explored and their ecological validity. In particular, only 6 studies relied on spontaneous speech data, while 4 other studies included a combination of spontaneous speech and (semi-)experimental approaches. Finally, Kašćelan and Parafita Couto identify lack of research in attitudes and advice that bilingual families with neurodivergent family members receive in terms of their code-switching practices. On rare occasions when this data is available, the advice and attitudes are predominantly negative, thus reinforcing stigmatisation of bilingual speech. 

Their chapter offers recommendations for future research, highlighting the importance of slow science for better understanding of code-switching by neurodivergent individuals. Importantly, in line with recommendations by \citet{Deuchar2020} among others, availability of spontaneous speech data is crucial in our shared understanding of this phenomenon, especially in cases such as neurodevelopmental conditions where the number of participants across existing studies has been particularly small. Making data available would allow for aggregated dataset creation, it would enable comparisons between different bilingual communities, and contribute to the validity of results. Advice and implications for educators and speech and language therapists are also offered based on the current findings in this field.

The contributions in this section collectively illustrate the intricate relationship between the characteristics of the languages, individual characteristics, caregiver strategies, and broader community practices in shaping bilinguals’ code-switching patterns.
\section{Concluding remarks}

The structure of this volume mirrors the intellectual trajectory and influence of Margaret Deuchar in the study of code-switching. Beginning with investigations into grammatical variation and constraints on code-switching, the volume reflects Deuchar’s pioneering contributions to the linguistic patterns governing bilingual speech (e.g., \citealt{Deuchar2006, DeucharEtAl2007}). Her work on Welsh-English, Spanish-English and Spanish-Welsh bilingualism and corpus-based methodologies laid the foundation for understanding how different languages interact at a structural level (\citealt{Bangortalknodate, Deuchar2020}).

Furthermore, this volume extends Deuchar’s legacy by emphasizing how code-switching operates in dynamic linguistic and social contexts. Deuchar’s more recent work has increasingly incorporated interdisciplinary perspectives, bridging psycholinguistics, sociolinguistics, and corpus linguistics (e.g., \citealt{Vaughan-EvansEtAl2020}). This is evident in the studies presented here, which explore code-switching as a product of both cognitive processes and social interaction, in line with the broader scope she championed.

Finally, the inclusion of research on neurodevelopmental conditions highlights an emerging frontier in bilingualism studies, one that aligns with Deuchar’s advocacy for empirical, data-driven research and her commitment to making bilingual linguistic data widely accessible. The emphasis on spontaneous speech corpora and ecological validity throughout the volume echoes her pioneering role in corpus creation and open data initiatives (e.g., \citealt{DeucharEtAl2014,DeucharEtAl2018}).

In sum, this volume is not just a collection of studies on code-switching; it is a reflection of the evolution of the field itself, shaped in no small part by Deuchar’s contributions. Her influence permeates the chapters, reinforcing the idea that code-switching is a window into the complexities of bilingualism/multilingualism, one that continues to reveal new insights, just as her work has done over the decades.

\sloppy\printbibliography[heading=subbibliography,notkeyword=this]
\end{document} 
